\chapter{Installing and the first launch}

\section{Installation}

\subsection{Official channels}

\begin{enumerate}
  \item \textbf{GitHub Releases} (recommended): open
        \url{https://github.com/LucasCoraca/mobilecad/releases/latest}, download
        \file{MobileCAD-arm64.apk} and tap it in the system file manager to install. The first install
        requires allowing “install unknown apps” for that source.
  \item \textbf{F-Droid}: the project ships an F-Droid build recipe (\file{fdroid/}); just search for
        MobileCAD in F-Droid. The F-Droid build updates through F-Droid itself and does not use in-app
        updates.
  \item \textbf{Chinese beta} (the companion build for this book): built from the \file{i18n/zh-hans} branch
        with the package name \file{app.mobilecad.i18n}, it can be installed \textbf{alongside} the official
        version and runs a Simplified Chinese interface. It uses a separate test signature and is not part of
        official updates.
\end{enumerate}

\begin{warn}[Signatures and upgrades]
The official APK and any repackaged third-party build have different signatures, so they \textbf{cannot be
installed over each other}. If you installed the official version first and then want another build,
uninstall it first (documents saved in storage are unaffected; in-app settings and the recent-files list are
lost).
\end{warn}

\subsection{What happens on the first launch}

On the first launch the app does three things, usually within a few tens of seconds (about 20--40 seconds on a
tablet):

\begin{enumerate}
  \item \textbf{Unpack the runtime}: the roughly 1 GB of FreeCAD, Python, Open CASCADE and Qt runtime
        libraries inside the APK is unpacked into the app's private directory;
  \item \textbf{Start the engine}: it loads more than 100 native libraries, initialises Python 3.14 and
        imports FreeCAD; the splash screen shows \ui{Starting the CAD engine…}. The first launch after an app
        upgrade also shows \ui{Unpacking the CAD engine…} first;
  \item \textbf{Create an empty document}: a new document named \ui{Unnamed} opens with an empty feature tree.
\end{enumerate}

Every later launch takes only 2--5 seconds. The unpacking happens once (it repeats after an app upgrade if the
runtime version has changed). If the engine fails to start, the splash screen shows \ui{The CAD engine failed
to start: …} in red and keeps the reason so that you can report it.

\begin{tip}[No storage permission needed]
MobileCAD \textbf{does not request storage permission}. Opening, saving, importing and exporting all go through
Android's Storage Access Framework (SAF): the app can read and write whichever directory you pick. Documents
are saved where you choose; the app's private directory holds only configuration and crash recovery copies.
\end{tip}

\section{Up and running in thirty seconds}

After launching you will see the interface below (Chinese UI):

\screenshot{ui-main-zh.jpg}{The MobileCAD main screen (Chinese UI). The feature tree is on the left, the ☰ menu and document information at the top, the 3D view in the middle, the element tabs (Part / Assembly / Drawing / Variables / More) at the bottom, and the NaviCube in the top-right corner of the view.}

\begin{enumerate}
  \item \textbf{Spin it}: drag in the view with one finger — the model rotates around its centre.
  \item \textbf{Zoom in}: pinch with two fingers to zoom, drag with two fingers to pan.
  \item \textbf{Change the direction}: tap the \ui{Front}/\ui{Top}/\ui{Right} face of the NaviCube in the
        top-right corner and the view animates to that direction; drag the cube itself to rotate freely.
  \item \textbf{Open the menu}: tap ☰ in the top-left corner to see six menu groups: File, Edit, View, Tools,
        Macro and Help.
  \item \textbf{Draw something}: tap \ui{Sketch} in the top toolstrip, choose \ui{New sketch} in the toolbar
        that unfolds, pick a datum plane and you are in sketch mode (Chapter 4 walks through it in full).
\end{enumerate}

\section{Gesture reference}

MobileCAD's gestures differ between modelling and sketching. This is where it feels most different from
desktop CAD, so it is worth learning first:

\begin{center}
\small
\begin{tabularx}{\linewidth}{@{}lXX@{}}
\toprule
\textbf{Gesture} & \textbf{Modelling} & \textbf{Sketching} \\
\midrule
One finger drag & Rotate the view & Draw (the point that counts is 50 dp above your finger) \\
Two fingers drag & Pan & Pan (a second finger landing within 300 ms cancels the stroke) \\
Two fingers pinch & Zoom & Zoom \\
Three fingers drag & Rotate around the centre of the three fingers (large reorientations) & Same \\
One finger tap & Select a face, edge or vertex & Select geometry or a constraint \\
Double-tap & Clear the selection & Clear the selection \\
Two fingers tap & Open the context menu (the desktop right-click) & Same \\
Long-press & Precision pointer: a crosshair appears 60 dp above your finger so that it does not hide the
target & Long-press a dimension label to edit its value \\
Two fingers long-press & Box select (drag a rectangle to select objects) & — \\
\bottomrule
\end{tabularx}
\end{center}

\begin{tip}[Why is the pointer above your finger?]
Your finger hides the point it touches, and that is the hardest interaction problem in mobile CAD. In drawing
and in long-press precision mode, MobileCAD puts the actual point some distance above your finger (50 dp when
drawing, 60 dp on a long-press) and shows a crosshair with snapping hints. Once you are used to it, the
precision matches a mouse.
\end{tip}

\section{Interface language}

The app \textbf{follows the system language} automatically: Simplified Chinese if the system is Chinese,
English otherwise. You can also pin a language by hand:

\begin{enumerate}
  \item open \menupath{Edit → Preferences} (or tap \ui{Preferences} on the toolstrip);
  \item choose \ui{Language (current: …)};
  \item pick \ui{System language}, \ui{English} or \ui{Chinese (Simplified)} in the list. The interface
        redraws immediately and the document and engine state are kept.
\end{enumerate}

\screenshot[0.8]{ui-language-zh.jpg}{The language chooser. The first entry, “System language (…)”, tells you which language Android is using; the interface refreshes as soon as you switch.}

\begin{tip}[Consistent terminology]
Command names, task panel fields and property names in MobileCAD come straight from FreeCAD's own terms (for
example \ui{Pad}, \ui{Pocket}, \ui{Fully constrained}). If you also use desktop FreeCAD, you will see the same
terminology on both.
\end{tip}

\section{Setting up your workspace}

Three more settings you will use often live in \menupath{Edit → Preferences}:

\begin{itemize}
  \item \textbf{Unit system}: switch FreeCAD's unit scheme (millimetre/metre, inch and so on). It affects how
        every value is displayed and entered;
  \item \textbf{Theme}: follow the system, light or dark. Dark suits long modelling sessions and OLED screens;
  \item \textbf{Check for updates daily}: when on, the app quietly checks GitHub Releases at startup; if there
        is a new version, an update dialog appears (release notes, download size, and the reminder “Android
        closes MobileCAD during installation, so save your work first”).
\end{itemize}

\screenshot[0.8]{ui-preferences-zh.jpg}{Preferences: unit system, theme, language and update checks.}
